\section{Relevantna literatura}
\label{sec:lit}

\textbf{Skupovi podataka.} Klasični javni skupovi za segmentaciju žila mrežnice su DRIVE~\cite{staal2004} (40 slika, $565\times584$), STARE~\cite{hoover2000} (20 slika) i CHASE\_DB1~\cite{fraz2012ens} (28 slika $999\times960$, 14 djece, oba oka). Zbog malog broja slika rezultati na njima jako ovise o podjeli i lako se precijene. FIVES~\cite{jin2022fives} donosi 800 slika rezolucije $2048\times2048$ s pikselnim anotacijama, ravnomjerno podijeljenih na četiri kategorije (normalno, AMD, dijabetička retinopatija, glaukom), te uz svaku sliku ocjenu kvalitete (osvjetljenje/boja, zamućenje, nizak kontrast). Autori predlažu podjelu 600/200 koja uvažava te ocjene. Za CHASE\_DB1 postoje dvije neovisne ručne anotacije; slaganje dvaju anotatora (Dice $\approx0{,}78$) gruba je gornja granica koju model realno može doseći~\cite{fraz2012ens}.

\textbf{Klasične metode.} Pregledni radovi~\cite{fraz2012survey,moccia2018} dijele nenadzirane pristupe na filtriranje usklađenim filtrima, analizu Hessiana (\emph{vesselness}), praćenje žila i morfološke operacije, a nadzirane na klasifikaciju piksela iz ručno izrađenih značajki. Gotovo svi rade na \textbf{zelenom kanalu}, koji ima najveći kontrast između žila i pozadine, uz lokalno izjednačavanje histograma (CLAHE~\cite{zuiderveld1994}). Te dvije predobrade ispitujemo u eksperimentu~2.

\textbf{Duboko učenje.} Liskowski i Krawiec~\cite{liskowski2016} pokazali su da konvolucijska mreža koja klasificira patcheve nadmašuje klasične metode. U-Net~\cite{ronneberger2015} (koder--dekoder s preskočnim vezama) postao je standardna arhitektura za biomedicinsku segmentaciju jer dobro radi s malo slika i čuva fine detalje kroz preskočne veze. Galdran i sur.~\cite{galdran2022} pokazuju da vrlo mali, pažljivo treniran U-Net postiže gotovo iste rezultate kao znatno složeniji modeli, ali i da uspješnost znatno pada kad se testira na skupu drukčijem od onog za treniranje (\emph{cross-dataset}). To nam je važno jer kombiniramo dva izvora s različitim kamerama i rezolucijama. Kao jaču varijantu koristit ćemo U-Net s koderom ResNet-34~\cite{he2016} predtreniranim na ImageNetu, što je u praksi čest i jednostavan način za bolje značajke uz malo podataka.

\textbf{Funkcije gubitka i evaluacija.} Žile zauzimaju samo mali dio slike, pa se uz binarnu unakrsnu entropiju često koristi Dice gubitak~\cite{milletari2016}, koji je manje osjetljiv na neuravnoteženost klasa. clDice~\cite{shit2021} dodatno kažnjava prekide tankih žila i tako čuva povezanost (topologiju) mreže žila. Literatura~\cite{moccia2018,galdran2022} upozorava da točnost (\emph{accuracy}) na ovom zadatku malo govori jer je većina piksela pozadina. Zato kao glavnu metriku uzimamo Dice, uz preciznost, odziv, specifičnost i PR-AUC, i to računamo unutar vidnog polja.
